\documentclass[11pt]{article}
\usepackage[margin=1in]{geometry}
\usepackage[utf8]{inputenc}
\usepackage[T1]{fontenc}
\usepackage{amsmath,amssymb}
\usepackage{graphicx}
\usepackage{booktabs}
\usepackage{array}
\usepackage{natbib}
\usepackage{hyperref}
\hypersetup{colorlinks=true,linkcolor=blue,citecolor=blue,urlcolor=blue}

\emergencystretch=3em
\newcommand{\TVD}{\mathrm{TVD}}
\newcommand{\pending}[1]{\textit{Results: pending (produced by \texttt{#1}).}}
\newcolumntype{L}[1]{>{\raggedright\arraybackslash}p{#1}}

\title{The Real World Benchmark: Forward-Dated Evaluation of\\Population Simulators Against Polls, Markets and Pricing Experiments}
\author{Real World\\\small Draft for the NeurIPS Datasets and Benchmarks Track. Not yet peer reviewed.}
\date{Draft of 11 October 2026 (second revision after internal review)}

\begin{document}
\maketitle

\begin{abstract}
Systems that simulate human populations with large language models (LLMs) are now marketed as complements or substitutes for surveys and market research \citep{brand2025llms,researchlive2026strat7}. Most evidence for them comes from static benchmarks whose answers are already public, and often already inside the training data of the models being tested. We pre-register the Real World Benchmark, a forward-dated benchmark with three tracks. Track~A asks a system to reproduce the published results of recurring opinion polls, including demographic subgroups, with every prediction committed before the wave's fieldwork begins. Track~B scores probabilistic calls on future events, generated by a fixed rule and committed into fixed one-per-cell Merkle slots, against a time-weighted prediction-market midpoint at a commit time attested by two independent RFC~3161 time-stamping authorities. Track~C asks for the outcome class and lift of subscription pricing and paywall experiments committed before the first user is exposed. Every score is reported next to cheap baselines: the naive last result, a single prompt to a frontier model and to the simulator's own backbone, a same-feed LLM given the simulator's data and number of calls, the market price where one exists, and, for Track~A, uniform and demographic-lookup anchors. We specify rounds, item-generation rules, paired score differences that are unbiased under pure sampling noise in the published estimate, seven confirmatory one-sided tests with Holm correction within each track, minimum detectable effects, and leaderboard rules against abstention, backfilling and selective reveal. No system has been scored. Every result is marked pending with the artifact that will produce it.
\end{abstract}

\section{Introduction}

A growing class of systems claims to predict what people will think, buy or do by simulating them. Silicon samples condition an LLM on a demographic profile and read off survey answers \citep{argyle2023outofone}. Digital twins condition on a person's past answers \citep{park2024agents,toubia2025twin}. Agent populations combine many such respondents into a society and let news move them. The company behind this paper is building one of the last kind (Section~\ref{sec:system}). As designed, it will hold about 100{,}000 synthetic agents as numeric records, built from the open synthetic demographic population GLOPOP-S \citep{ton2024glopop} and published survey aggregates rather than from census microdata, whose terms bar commercial use \citep{ipumsi2026}; it will correct segment beliefs each morning against the polls released that day, with news as forcing, and forecast from that state. That simulator does not exist yet; the hypotheses below that involve it can be tested only if it is registered before round~1 opens.

Evidence for all of these systems is weak in a specific way. The main benchmarks score static distributions from survey waves published years ago \citep{santurkar2023whose,durmus2023global,hu2026simbench}. The models being scored have plausibly read those waves, so a good score mixes simulation with recall \citep{sainz2023nlp}, and declared knowledge cutoffs need not match effective ones \citep{cheng2024dated}. The benchmarks also rarely publish what a cheap non-LLM predictor scores on the same items. When that comparison is made, the picture darkens: in a test of four models from two families on General Social Survey and World Values Survey items, no LLM beat the strongest fitted non-LLM baseline at the individual level, and the models inflated between-segment gaps two- to fourfold \citep{chen2026synthetic}. The best of 45 models on SimBench scores 40.80 out of 100, where 0 is a uniform guess \citep{hu2026simbench}. A trade-press report of an industry white paper, not peer reviewed, describes synthetic willingness-to-pay about 16\% above human answers and wave-to-wave change tracked correctly 47\% of the time \citep{researchlive2026strat7}.

Forecasting research has a better habit. ForecastBench asks about events unknown at submission time \citep{karger2024forecastbench}; Prophet Arena scores LLMs against the Kalshi price \citep{yang2025prophet}; Metaculus runs bot tournaments against professionals \citep{wilson2026futureeval}. These score a probability on an event, not a distribution of opinions across subgroups, and not the outcome of a business experiment.

The Real World Benchmark joins the two traditions. It keeps the rule that a prediction counts only if it was made before the answer existed, and applies it to the outputs population simulators sell: subgroup opinion distributions, event probabilities and pricing outcomes.

\paragraph{Contributions.}
\begin{enumerate}
  \item \textbf{Three forward-dated tracks} with item schemas, item-generation rules, per-source commit deadlines and void rules: Track~A (recurring polls with subgroups, committed before fieldwork), Track~B (rule-generated event questions against markets), Track~C (pricing and paywall experiment outcomes, committed before first exposure).
  \item \textbf{A fixed baseline set}, including a single prompt to the simulator's own backbone and, in Tracks~A and~B, a same-feed LLM given the same observations, history window and number of calls, so a simulator's state and population are compared with the strongest cheap rival while data and compute are held fixed.
  \item \textbf{Scores built for noisy truth}: a multiclass squared error for Track~A whose paired differences are unbiased for the difference in squared distance to the poll's expected value under pure sampling noise; a noise-corrected subgroup gap error whose paired differences are unbiased in the same sense; paired Brier differences with restricted wild cluster bootstrap tests; and ranked probability scores for Track~C, which are proper but not corrected for noisy truth.
  \item \textbf{Leaderboard rules} that close escape routes: one leaf per (entry, item, horizon) at a fixed slot, 128-bit salts, unrevealed commitments scored as the worst allowed outcome, intention-to-treat imputation that never rewards abstention, an external randomness beacon for item sampling, a hashed system card with an information cutoff covering every input, separate divisions and modes, and fixed analysis dates.
  \item \textbf{A pre-registered analysis plan}: seven confirmatory one-sided tests, Holm-corrected within each track at a pre-split level, with stated minimum detectable effects, a list of every parameter still open and when it is frozen, estimation targets reported without tests, and the artifact path that will hold each result.
\end{enumerate}

\paragraph{Scope and companion papers.} This paper defines the benchmark. Companion papers in preparation (short titles in italics) cover the simulator as a whole and its matched comparison with asking a model (\emph{Don't Ask a Model, Ask a World}), the commitment and timestamping protocol for its own daily record (\emph{Receipts}), an upgraded filter that moves beliefs toward observations (\emph{Belief Assimilation}), the population resource (\emph{Unequally Known}), corrections for narrow spread (\emph{Correcting Variance Collapse}), a survey of silicon samples (\emph{Silicon Samples, 2022--2026}) and the governance of public world simulators (\emph{Governing Public World Simulators}). The rules in Section~\ref{sec:rules} are binding on their own and do not depend on those papers.

\section{Related work}

\paragraph{Opinion simulation benchmarks.} OpinionQA measured how far LM answer distributions sit from 60 US demographic groups on Pew American Trends Panel questions \citep{santurkar2023whose}. GlobalOpinionQA extended this to cross-national Pew and World Values Survey items and found that default model answers resemble those of the USA and some European and South American countries more than others \citep{durmus2023global}. WorldValuesBench builds a large WVS-derived dataset \citep{zhao2024worldvaluesbench}. SubPOP fine-tunes on survey data to predict subpopulation distributions \citep{suh2025subpop}. SimBench unifies 20 datasets, scores 45 LLMs with a normalized total-variation score, and reports that simulating specific demographic groups is harder than simulating the general population \citep{hu2026simbench}. Twin-2K-500 and interview-grounded agents test whether a model can stand in for a specific person \citep{toubia2025twin,park2024agents}. The Artificial Societies Benchmark is the closest in aim to Track~A: it combines eleven validity tests drawing on twenty human sources across nine models, and finds that fidelity in one domain does not establish fidelity in others, with models compressing response scales and altering relationships between traits \citep{chidichimo2026artificial}. All of these score answers that were public before the models were scored, most before they were trained. Track~A differs in one respect only, but it is the one that matters for contamination: its truth does not exist when predictions are committed. It reports the SimBench score as a description and tests hypotheses with a different score (Section~\ref{sec:metrics-a}).

\paragraph{Critiques of silicon samples.} Synthetic ANES respondents reproduce average feeling-thermometer scores but vary too little and give different regression results \citep{bisbee2024synthetic}. LLM survey answers shift with answer order and, once that is controlled, many models give near-uniform answers \citep{dominguezolmedo2023questioning}. LLMs asked to stand in for identity groups misportray and flatten them \citep{wang2025flatten}, and an evaluation of digital twins organizes their failures into five distortions \citep{peng2025funhouse}. Chen et al.\ find that no LLM beats the strongest fitted baseline at the individual level and that models inflate between-segment differences \citep{chen2026synthetic}; their primary yardstick is a demographic lookup with back-off, which our Track~A anchor follows. These results motivate three design choices: every score ships with a non-LLM baseline, subgroup gaps are scored directly, and the single-prompt baseline averages over answer orders.

\paragraph{Forecasting benchmarks.} Retrieval-augmented LLM forecasters approach the crowd aggregate of competitive forecasters \citep{halawi2024approaching}. ForecastBench found that expert forecasters outperformed the top LLM \citep{karger2024forecastbench}; its maintainers reported in January 2026 that superforecasters still led the best LLMs by 0.017 Brier points and projected parity around November 2026, with a 95\% interval from January 2026 to November 2027 \citep{fri2026closing}. Prophet Arena reports, over 1{,}367 resolved Kalshi events, a best-model Brier score of 0.184 against 0.187 for the market, each $\pm 0.006$ at 95\% \citep{yang2025prophet}. In the Metaculus FutureEval spring 2026 tournament (173 bots, 297 scored questions), the best baseline bot, which used AskNews retrieval and aggregated five forecasts, placed 18th \citep{wilson2026futureeval}. Earlier event-forecasting datasets include Autocast \citep{zou2022autocast}. Paleka et al.\ catalogue how leakage and logical shortcuts inflate forecaster evaluations \citep{paleka2025pitfalls}, which informs our commit and eligibility rules.

\paragraph{Markets as a reference.} Prediction-market prices are widely treated as a strong aggregate forecast \citep{wolfers2004prediction}, but a price need not equal traders' mean belief \citep{manski2006interpreting}, and prices show a favourite-longshot pattern \citep{snowberg2010,page2013calibrated}. We use the market as a hard bar, not as ground truth, and test sensitivity to the liquidity rule and the reference window.

\paragraph{Predicting experiment results.} Ashokkumar et al.\ prompted GPT-4 to simulate respondents in 70 preregistered US survey experiments; predicted effects correlated strongly with actual effects, also for studies not public at the training cutoff, but systematically overestimated effect sizes \citep{ashokkumar2026predict}. DellaVigna and Pope compared expert and non-expert forecasts of experimental results \citep{dellavigna2018predicting}, and DellaVigna, Pope and Vivalt argue for collecting such forecasts before results are known \citep{dellavigna2019predict}. LLM conjoint responses can reproduce willingness-to-pay for existing features \citep{brand2025llms}. We know of no benchmark that scores predictions of real pricing or paywall experiments committed before the experiment starts; Track~C fills that narrower gap.

\paragraph{Dynamic benchmarks, timestamping and pre-registration.} Dynabench and LiveBench refresh their items \citep{kiela2021dynabench,white2025livebench}. Our tracks are stricter: the answer does not exist at commit time. Hash-linked timestamping goes back to Haber and Stornetta \citep{haber1991timestamp}. We take minute-level commit times from RFC~3161 time-stamp tokens \citep{rfc3161}, which rest on trusting the issuing authority, and add OpenTimestamps \citep{opentimestamps}, whose Bitcoin attestations need no trusted party but are accurate only to hours \citep{todd2016ots}. We follow pre-registration practice, including a dated public log of deviations \citep{nosek2018preregistration}.

\section{Design principles}

\begin{enumerate}
  \item \textbf{The answer must not exist when the prediction is committed.} Only live items enter tests and ranks; any live item whose answer turns out to have existed at its commit time is void. Retrospective items are reported in separate columns for development.
  \item \textbf{Items are generated by rule, not chosen.} The organizer also enters, so item selection follows published code hashed before the round opens, and random draws use an external beacon.
  \item \textbf{Baselines first.} A score is reported only next to the baselines of Section~\ref{sec:baselines}, on the same items.
  \item \textbf{Subgroups are scored where subgroup truth exists, and only there.}
  \item \textbf{Every committed prediction counts.} Misses, missing and unrevealed predictions are all scored, in tests as well as ranks.
  \item \textbf{Inference happens once per track.} Tests and intervals are computed on a fixed analysis date per track and round, never on interim data.
\end{enumerate}

\section{Rounds, roles and system cards}
\label{sec:rounds}

\paragraph{Rounds.} Round $r$ is defined by a manifest hashed and timestamped at least 14 days before its open time $O_r$. The manifest fixes every parameter in Table~\ref{tab:free}. Round~1 is planned to open on 1 December 2026; that date is not yet confirmed.
\begin{itemize}
  \item \emph{Track~B}: items are generated weekly for 26 weeks from $O_r$; the analysis date is $O_r + 32$ weeks. Items unresolved by then are void.
  \item \emph{Track~A}: covers waves whose commit deadline (Section~\ref{sec:track-a}) falls within 26 weeks of $O_r$; the analysis date is $O_r + 60$ weeks. Pew releases one spring wave in stages: the 2026 wave, whose fieldwork ended on 13 May 2026, appeared in reports from 23 June to 17 September 2026 \citep{pew2026methods,pew2026ai}. Sixty weeks from a 1 December open gives about eight months after a similar 2027 fieldwork end. Items unpublished by then are void, and void rates are reported by topic, because late release may correlate with topic.
  \item \emph{Track~C}: live calls committed within 52 weeks of $O_r$; the analysis date is $O_r + 65$ weeks.
\end{itemize}
Void items are not scored for anyone and are counted in the report. If a minimum count (Section~\ref{sec:analysis}) is not reached on the analysis date, the test is reported as not tested; it is never run early, and a later round may pre-register a pooled test only by amendment before its own open date.

\paragraph{Roles and information barrier.} The \emph{operator} runs item generation, the commit log, market snapshots and scoring from a repository the simulator team cannot write to. Items are released to every entrant, the organizer's simulator team included, through one public feed at the same moment. Because the operator and the simulator team belong to one company, this barrier is organizational, not independent (Section~\ref{sec:limits}). Until an external entrant registers, a round is a self-evaluation of the organizer's system against baselines and is labelled as such, not as a leaderboard.

\paragraph{System card and information cutoff.} Before $O_r$, each entry registers a system card whose hash is published: organization, model components and versions, code commit and configuration hash, retrieval sources, population source data, assimilation inputs, whether market prices are read, ensemble size and number of model calls per item, cost per item, and any positions its organization holds in markets on benchmark items. The entry's \emph{information cutoff} is the latest date of any information the system uses: each model's declared training cutoff, the retrieval-index snapshot, the latest fieldwork in its population source data, the latest assimilation input, and the latest data used for tuning or hyperparameter selection. For live items the commit deadline bounds every dynamic input; the information cutoff governs retrospective eligibility.

\paragraph{Deviations log.} Every change to this protocol after publication is a dated, hashed entry in an amendments file, stating what changed, why, and which rounds it applies to. No amendment applies to a round already open.

\paragraph{Freezing this text.} This draft is not yet hashed. We plan to hash it and obtain RFC~3161 and OpenTimestamps attestations at the commit that freezes it, before the round-1 manifest and before the first resolved question of the organizer's own daily record. Table~\ref{tab:free} lists every parameter this text leaves open and where it will be fixed; none may change after that point except through the deviations log.

\begin{table}[t]
\centering
\small
\begin{tabular}{@{}L{6.4cm}L{4.0cm}L{3.6cm}@{}}
\toprule
\textbf{Open parameter} & \textbf{Fixed in} & \textbf{Frozen by} \\
\midrule
Round-1 open date $O_1$ (planned 1 Dec 2026) & Round manifest & Manifest hash, $O_1 - 14$ days \\
Ipsos commit schedule (rule in Section~\ref{sec:track-a}) & Round manifest & $O_1 - 14$ days \\
Track~B calendar list (central banks, statistics series), market categories, $K$ per week & Round manifest & $O_1 - 14$ days \\
Time-stamping authorities, drand chain and round offset & Round manifest & $O_1 - 14$ days \\
Frontier and backbone model versions, successors & Round manifest & $O_1 - 14$ days \\
Prompt texts T1--T3; same-feed LLM token budget and history window & Round manifest & $O_1 - 14$ days \\
RW-blind keyword filter; organizer ensemble size $E$ & Organizer system cards & $O_1$ \\
\bottomrule
\end{tabular}
\caption{Every parameter this paper leaves open, where it is fixed and by when.}
\label{tab:free}
\end{table}

\section{The three tracks}

Table~\ref{tab:tracks} summarizes the tracks. All items share an envelope: item id, track, as-of time $\tau$ (the latest information a prediction may use, which is also the commit deadline), resolution source, and a content hash of the item frozen when it is released.

\begin{table}[t]
\centering
\small
\begin{tabular}{@{}L{2.1cm}L{4.1cm}L{4.1cm}L{4.1cm}@{}}
\toprule
 & \textbf{A: Poll replication} & \textbf{B: Live calls vs markets} & \textbf{C: Pricing decisions} \\
\midrule
Predict & Distribution per (question, country, subgroup) & $P(\text{event})$ at 28, 7 and 1 days & Outcome class, lift bucket, metric directions (all as probabilities) \\
Truth & Published estimates of the next wave of the same poll & Resolution source named in the item & Partner-reported result \\
Commit by & Per-source date before fieldwork can start & Snapshot time $\tau$ per horizon & Before first exposure \\
Primary score & Multiclass squared error $Q$ & Paired Brier difference & Ranked probability score \\
Baselines & Last wave, single prompt (frontier, backbone), same-feed LLM, uniform, demographic lookup & Last result, base rate, single prompt, same-feed LLM, market & Running prior, single prompt \\
\bottomrule
\end{tabular}
\caption{The three tracks of the Real World Benchmark.}
\label{tab:tracks}
\end{table}

\subsection{Track A: recurring polls with subgroups}
\label{sec:track-a}

\paragraph{Task and estimand.} Given a survey question, its options, a country, a subgroup (for example ``ages 18 to 34'') and $\tau$, output a probability distribution over the options. The estimand is what the poll will report, not the population's true opinion: Ipsos samples in 14 of its countries are more urban, more educated or more affluent than the general population \citep{ipsos2026may}, and predictions are scored against those samples, in each poll's own age frame and country frame as stated in each wave's methodology.

\paragraph{Eligible items.} Live items are \emph{trend items}: every eligible question with a published distribution in wave $w$ of a named recurring series becomes a candidate for wave $w{+}1$ in every country of wave $w$. Eligible are single-choice binary, nominal and ordinal questions. A multi-select question (``which three of the following\ldots'') becomes one binary row per option, \emph{selected} or \emph{not selected}, because its shares sum to more than one and do not form a distribution. Numeric and open questions are excluded. If the source publishes a ``don't know'' or refused share for an item, it is an option to predict. If it does not, $h$ is the published substantive shares renormalized to sum to one; such rows are flagged, because dividing by a noisy total makes the unbiasedness of Section~\ref{sec:metrics-a} approximate. A single-choice row whose published shares sum outside $[0.97, 1.03]$ without a stated exclusion is void.

\paragraph{Round-1 sources and commit deadlines.}
\begin{itemize}
  \item \emph{Pew Research Center Global Attitudes spring survey} (annual). Recent waves were fielded 14 February--3 June 2022 (19 countries), 20 February--22 May 2023 (23), 5 January--22 May 2024 (35) and 8 January--26 April 2025 (25) \citep{pewdatasets2026}; the 2026 wave surveyed 42{,}151 people in 36 countries from 8 February to 13 May 2026 (37 when the separately fielded US surveys are counted), with two separate US surveys (5{,}119 adults, 17--23 February; 3{,}488 adults, 22--28 June) \citep{pew2026ai}, and publishes a per-country fieldwork timeline \citep{pew2026methods}. Reports carry subgroup splits, but not uniformly: the AI report discusses gender where the gap is significant (in 11 mostly high-income countries women were more likely than men to expect fewer jobs) and ideology in the 28 countries where it is measured \citep{pew2026ai}. The commit deadline for the spring 2027 wave is $\tau = $ 1 January 2027, 00:00 UTC, four days before the earliest start of the last four waves.
  \item \emph{Ipsos What Worries the World} (monthly, 30 countries, about 20{,}000 to 24{,}000 adults aged 16--74 per wave, online panel) \citep{ipsoswww2026,ipsos2026may}. Fieldwork for a wave can start in the month before the one it is named after: the May 2026 wave was fielded 24 April--8 May 2026 \citep{ipsos2026may}. The commit deadline for each wave is 14 days before the earliest fieldwork start, relative to the wave's named month, documented across the twelve waves before $O_r$; the manifest lists the resulting dates. Its main question is multi-select (``which three'') and enters as per-option binary rows; the ``right direction or wrong track'' question is binary. Country-by-issue tables are published; demographic breakdowns are available only on request, so Ipsos contributes no subgroup rows.
\end{itemize}
\emph{Void rules.} A (wave, country) whose documented fieldwork start precedes $\tau$ is void; this applies to every source, and the Pew deadline is an extrapolation that this rule backstops. An item is also void if the question is not repeated with identical wording and options, the country is not fielded again (Pew's country set changes from year to year), or no report carrying the item is published by the analysis date. An item's publication time is that of the first report or topline that carries its distribution for that country and subgroup. Commitments on void items are disclosed but not scored. The yield is not known in advance; a feasibility count of Pew 2025 questions repeated in 2026, with the share of (question, country) cells lost to country turnover, is a pre-submission artifact (Table~\ref{tab:pending}), and the realized yield is a round-open census. Subgroup items exist only for the Pew source in round~1, so if fewer question clusters with subgroup truth survive than the minimum in Section~\ref{sec:analysis}, the subgroup hypotheses are not tested.

\paragraph{Row schema.} One row per (source, wave, question, country, subgroup, option row for multi-select items): question text and options in the fielded and English versions, item type, fieldwork dates, publication date, frame, reported $n$ and margin or interval where published, the published distribution $h$, a subgroup-truth flag, the demographic dimension of the subgroup, source URL, retrieval time and a hash of the raw source file. Subgroup rows with reported $n < 100$ are excluded. Rows whose redistribution terms are unclear are released as links and hashes.

\paragraph{Retro mode.} A published poll is eligible for an entry in retro mode only if its fieldwork started after the entry's information cutoff, with retrieval and assimilation frozen at that cutoff. Retro scores are reported in a separate column and never enter tests or ranks. The Pew spring 2026 wave started fieldwork on 8 February 2026; it is ineligible for any entry whose information cutoff is on or after that date, which includes the organizer's simulator, and serves only as a development set.

\subsection{Track B: rule-generated calls against markets}

\paragraph{Task.} Given a binary question with a fixed resolution rule and source, output a probability at 28, 7 and 1 days before scheduled resolution. Submissions must lie in $[0.01, 0.99]$. Scheduled resolution is the meeting or release time for a calendar item and the contract's scheduled close or expiration time, as stated in its rules at generation, for a market item.

\paragraph{Item generation.} Items come from two rules in the hashed generation code, run weekly; each week's candidate list is published and timestamped before it is sampled.
\begin{itemize}
  \item \emph{Calendar items}: scheduled events from official calendars named in the manifest, instantiated from fixed template types. Each type defines its threshold, its status-quo outcome and its base rate:
  \begin{itemize}
    \item \emph{Rate decision}: ``Will central bank $C$ change its policy rate at the meeting on $D$?'' No threshold. Status quo: no change. Base rate: the share of $C$'s scheduled meetings in the five years before $O_r$ with no change.
    \item \emph{Release direction}: ``Will the first release of statistic $S$ for period $P$ exceed the first release for the preceding period?'' The threshold is that earlier first release, which is published before generation; a tie resolves NO. Status quo: the direction of the most recent change repeats. Base rate: the share of first releases of $S$ in the five years before $O_r$ that exceeded the previous first release.
  \end{itemize}
  A new template type enters only through the manifest, with the same three fields.
  \item \emph{Market items}: every binary Kalshi or Polymarket contract in the manifest's named categories that, at the weekly generation time, resolves 35 to 42 days later, names a resolution source and meets the coverage rule below. Contracts are grouped by the venue's event identifier, so a strike ladder or a multi-outcome event is one group. If more than $K$ groups qualify in a week, $K$ groups are drawn uniformly, and one contract is drawn uniformly within each sampled group. The seed is the drand beacon output \citep{drand} for the first round at least one hour after the candidate list is published; this applies from week~1 and does not depend on any value the organizer can influence. The item's resolution rule is the market's own rules text.
\end{itemize}
Generation 35 to 42 days ahead leaves at least seven days between public release and the 28-day snapshot. A calendar item is matched to a market only if the contract names the same event, threshold and resolution source. Items are clustered by rule: items sharing a meeting, release, election or market event identifier form one cluster.

\paragraph{Void rules.} A horizon slot is void if its outcome was determined before $\tau_{i,h}$: the event occurred, the contract closed or settled early, or the official figure was published. Items whose market resolves 50-50, ``too early'', or remains disputed on the analysis date are void. The number of voided slots is reported per horizon and rule.

\paragraph{Commitment.} The snapshot time $\tau_{i,h}$ is 12:00~UTC on the day $h$ days before scheduled resolution. The binding rules, for every entry including the organizer's:
\begin{enumerate}
  \item The operator builds the tree. Each registered entry has exactly one slot per (item, horizon), at a position fixed by sorting (entry id, item id, horizon) in the manifest.
  \item An entrant submits a leaf $H(\text{canonical}(\text{entry}, \text{item}, h, p) \,\|\, \text{salt})$ before $\tau$, where canonical is the JSON Canonicalization Scheme \citep{rfc8785} and the salt is at least 128 bits from a cryptographically secure generator, so that the few plausible values of $p$ cannot be recovered by brute force. The operator's append-only log publishes each received leaf hash with its receipt time at once and rejects a second leaf for an occupied slot.
  \item At $\tau$ the operator places in each empty slot a filler leaf, $H(\texttt{"empty"} \,\|\, \text{canonical}(\text{entry}, \text{item}, h))$. It then publishes the batch root and its leaf count, which equals the number of slots due, and requests RFC~3161 time-stamp tokens \citep{rfc3161} on the root from at least two independent time-stamping authorities named in the manifest, none operated by the organizer. The provable commit time $t^\ast$ is the later of the two token times, so backdating requires both authorities to collude. The root is also submitted to OpenTimestamps \citep{opentimestamps}. Its Bitcoin attestation usually arrives hours later and the block time is accurate only to a few hours \citep{todd2016ots}; it serves as a trust-free upper bound, and a batch whose attesting block time is more than six hours after $\tau$ is flagged in the report.
  \item If $t^\ast \le \tau + 10$ minutes, the market reference is taken at $\tau$. Otherwise each prediction in the batch is scored against whichever reference, at $\tau$ or at $t^\ast$, gives the market the lower Brier score.
  \item After resolution the entrant reveals $(p, \text{salt})$. A committed leaf that is never revealed, or whose reveal does not match, is scored as the worst allowed submission, $p = 0.99$ on the wrong outcome, for every score. An empty slot is a missing prediction (rule~\ref{rule:missing} in Section~\ref{sec:rules}).
\end{enumerate}
The organizer's own daily public record follows the \emph{Receipts} companion, which reads its market reference from venue history at a scheduled call time; benchmark scores use only the rules here, and \emph{Receipts} may strengthen them for the organizer's entries but not weaken them.

\paragraph{Market reference and coverage.} The operator snapshots each order book every minute, hashing each snapshot into a sibling tree of the batch. The Kalshi order-book endpoint returns only yes bids and no bids; the yes ask is derived as one minus the best no bid \citep{kalshi2026orderbook}. Polymarket books come from its central limit order book \citep{polymarket2026books}. The reference $q_{i,h}$ is the time-weighted bid-ask midpoint over the 60 minutes before $\tau$; because entrants may use information up to $\tau$, a sensitivity analysis uses the last 15 minutes. An item is \emph{market-covered} at a horizon if its time-weighted spread is at most 0.05 and at least US\$1{,}000 traded in the preceding 24 hours. A sensitivity analysis uses a spread of at most 0.10. Uncovered items are scored against the other baselines and reported separately; because liquidity correlates with topic and difficulty, covered and uncovered scores are not compared with each other.

\paragraph{Divisions.} The market-blind division forbids prediction-market prices in any input. The market-aware division allows them. The organizer registers two configurations of its simulator, each with its own hashed system card:
\begin{itemize}
  \item \emph{RW-blind}, in the market-blind division: no prediction-market price in any input or in assimilation, and news items that mention a listed market platform or betting odds dropped by a published keyword filter.
  \item \emph{RW-aware}, in the market-aware division: the same system with market prices allowed as inputs.
\end{itemize}
All Track~B hypotheses about the simulator use RW-blind. In the market-aware division, matching the market is trivial by copying it.

\subsection{Track C: pricing decisions (Paywall Bench)}

\paragraph{Task.} Given a blinded description of a subscription app's pricing or paywall experiment (control and variant arms, what was held constant, the primary metric and its better direction, secondary metrics, the planned sample size or duration, the minimum detectable effect the partner planned for, and the partner's pre-declared significance level $\alpha$, default 0.05 two-sided), output:
\begin{itemize}
  \item probabilities over three ordered outcome classes on the primary metric: control better, no detectable difference, variant better, where ``better'' means the partner's interval at level $\alpha$ excludes zero in the metric's declared better direction (a decrease for lower-is-better metrics such as churn);
  \item probabilities over six ordered buckets of the primary metric's raw relative change, variant against control: $<-10\%$, $-10$--$0\%$, $0$--$10\%$, $10$--$30\%$, $30$--$100\%$, $\ge 100\%$;
  \item for each secondary metric, probabilities over down, flat and up, where flat means the partner's interval includes zero, or, if no interval is reported, a relative change below 2\% in absolute value.
\end{itemize}
The outcome class depends on the test's power as well as on the effect, which is why the planned sample size and minimum detectable effect are part of the item.

\paragraph{Truth.} The partner reports the point estimate and interval for each metric, the stopping date, and whether the test stopped at its pre-registered end. The ship decision is recorded but not scored. Outcomes are self-reported and cannot be verified by third parties.

\paragraph{Live set.} The partner writes the blinded description; the operator only removes identifying details. The description is hashed and frozen before the first user is exposed to either arm, and calls are committed by hash, under the Track~B slot and salt rules, before that same time. Before submitting a test, a partner registers its planned end date or sample size and agrees to publication of the blinded description, outcome class and lift bucket. Tests whose partners do not agree are not admitted, and a test whose first exposure precedes the commit time is void. Partners cannot read predictions before the test ends, because only salted hashes are public. Tests stopped before their registered end are flagged and reported separately. No user-level data or revenue figures are stored.

\paragraph{Retrospective set.} An internal pilot of public A/B tests, reconstructed from pages published by app makers and subscription platforms, exists in a separate repository. No counts from it are reported here until it is imported into the benchmark repository and each row re-verified. Retro items are split per entry: an item is \emph{post-cutoff} for an entry if the source page was published after the entry's information cutoff, otherwise \emph{pre-cutoff}. Published case studies are unlikely to be representative: in Microsoft's experimentation platform, roughly a third of ideas were positive and significant, a third flat and a third negative \citep{kohavi2009tutorial,kohavi2020trustworthy}, while app makers and vendors choose which of their tests to publish. The retrospective set therefore supports no inference and is used only for development.

\section{Baselines and the organizer's entry}
\label{sec:baselines}

Each baseline is a frozen program whose code and prompts are hashed in the round manifest. Baselines are unranked and run at $\tau$ like entrants.

\paragraph{The organizer's entry.}
\label{sec:system}
As designed in \emph{Don't Ask a Model, Ask a World}, the organizer's system holds about $10^5$ agents as numeric records: a demographic record from GLOPOP-S \citep{ton2024glopop} and a fixed offset per modelled topic. Agents are grouped into up to about 1{,}800 segments (country $\times$ age band $\times$ education). The state is one latent position per segment on each topic in a fixed list, plus pollster house effects. Each day an LLM grades each news story once into a signed signal that moves segment positions, and a Kalman-type filter corrects the positions against newly released polls; prediction-market prices are never assimilated. Predictions are made in three ways, and the share of items made each way is reported:
\begin{itemize}
  \item A poll row on a modelled topic is computed without an LLM call, by applying a response model to the agents in the poll's frame. This is where the agent records enter: they represent the spread of opinion within each segment.
  \item An event that the population decides (a vote share or a poll threshold) gets the share of $E$ ensemble members whose computed share falls on the YES side.
  \item Every other event, and every poll question outside the topic list, gets one backbone call per ensemble member, given the item, the observation bundle and a summary of the state; the prediction is the mean over the $E$ members.
\end{itemize}
The model is never called per agent. ``Population'' in the hypotheses below therefore means the state and its agent records together, not a set of simulated respondents. The engine built so far (version 0.1.0) is simpler than this design: $10^4$ agents sampled from illustrative marginals written for its repository, a scalar Kalman update of each country's mean belief per topic, and one backbone call per segment for every item, with no response-model path. The system card at $O_1$ describes whichever version is entered.

\paragraph{Naive last result.}
\begin{itemize}
  \item \emph{Track A:} the distribution from the most recent earlier wave of the same question, country and subgroup published before $\tau$; if the subgroup is missing, that wave's country topline; if the question is new, uniform. Fallback coverage is reported.
  \item \emph{Track B (calendar items only):} probability 0.8 on the template's status-quo outcome, with 0.7 and 0.9 as sensitivity checks. It is undefined for market items.
  \item \emph{Track C:} the running prior: Laplace-smoothed frequencies of the three outcome classes and six lift buckets among live tests whose outcomes were recorded in the ledger before this item's commit time; uniform when there are none. The publication-biased retrospective pool is not used.
\end{itemize}

\paragraph{Base rate (Track B calendar items).} The template's five-year base rate defined above, placed on the status-quo outcome. It is the comparator for H6.

\paragraph{Exploratory cheap baseline.} Track~A \emph{last wave plus drift}: the last wave plus the cross-country mean change for the same question between the two previous waves, projected onto the simplex.

\paragraph{Single prompt.} One fixed template, T1, run on two models: the frontier model (\emph{SP-frontier}) and the simulator's own backbone (\emph{SP-backbone}), each named with its exact version in the manifest. The frontier model is, among models generally available through a public API 14 days before $O_r$, the one with the latest documented knowledge cutoff, ties broken by the later general-availability date. The prompt contains only the item's information packet: no tools, no retrieval, and no market price in the market-blind division. Elicitation is verbalized: the model returns a probability (Track~B) or a distribution (Tracks~A and~C) as JSON at temperature~0, or at the provider's fixed setting where temperature cannot be set. One sample is drawn with options in source order and one with options reversed, and the two are averaged \citep{dominguezolmedo2023questioning}. Invalid output is retried once, then imputed as for a missing prediction. Two further templates, T2 and T3, are run and their spread reported; hypotheses use T1 only. If a named model is withdrawn mid-round, the manifest's designated successor takes over, and results are reported pooled and split by version.

\paragraph{Same-feed LLM (Tracks A and B).} The simulator's backbone, called $E$ times, where $E$ is the organizer's ensemble size from its system card, with the same aggregation rule (the mean of the returned distributions or probabilities). Each call receives the observations the organizer's entry received over the same history window: for Track~A, every poll aggregate it assimilated up to $\tau$, including the latest published wave of the item's question, and the news it used as forcing; for Track~B, RW-blind's assimilation inputs. Inputs are ranked by the simulator's own relevance score and truncated to a token budget fixed in the manifest. There is no state and no population. Data, history window, retrieval ranking and number of calls are thus held fixed. One difference remains by construction: the simulator integrates the whole history in its state, while the LLM sees only what fits in the budget. Comparisons with this baseline therefore test the state, the population and that integration together.

\paragraph{Market price.} Track~B only: $q_{i,h}$ on market-covered items.

\paragraph{Anchors for Track A.} The uniform distribution (for a per-option binary row of a choose-$k$-of-$K$ question, a selection share of $k/K$), and a demographic lookup with back-off in the manner of Chen et al.\ \citep{chen2026synthetic}, fit privately on the most recent respondent-level file containing the same question whose license permits this use. Cells are country $\times$ age band $\times$ sex $\times$ education with at least 50 respondents, collapsing to country $\times$ age band and then to the country topline, with age bands cut to each poll's frame. If no such file exists, the lookup has no value on that item and its coverage is reported. Fitted tables are released only where the license allows; otherwise code and hashes are released.

\section{Metrics and statistics}

\subsection{Track A}
\label{sec:metrics-a}

\paragraph{Primary score.} For row $i$ with published distribution $h_i$ over $K_i$ options and prediction $m_i$, the primary score is the multiclass squared error
\begin{equation}
Q_i = \textstyle\sum_k (m_{ik} - h_{ik})^2 ,
\end{equation}
which for a binary row is $2(m_i - h_i)^2$ in the share of the first option. Let $\bar h_i = \mathbb{E}\,h_i$ be the poll's expected value. Then $\mathbb{E}\,Q_i = \|\bar h_i - m_i\|^2 + \mathrm{tr}\,\mathrm{Cov}(h_i)$, and the second term does not depend on $m_i$. Paired differences in $Q$ between two predictors are therefore unbiased for their difference in squared distance to $\bar h_i$, however large the sampling noise. $\bar h_i$ is what the poll measures in expectation; it equals the population distribution only without house, frame, mode and wording effects, which we do not correct. $Q$ is proper for eliciting $\bar h_i$, so reporting one's expectation of the published result is optimal. Two departures are bounded rather than removed. Renormalization over substantive options (Section~\ref{sec:track-a}) makes $h_i$ a ratio and the unbiasedness approximate. Rounding to whole percentage points adds an error $r_{ik}$ with $|r_{ik}| \le 0.005$, which shifts a paired difference by $-2\sum_k (m_{1k} - m_{2k})\, r_{ik}$, at most $0.01 \sum_k |m_{1k} - m_{2k}|$ in absolute value.

\paragraph{Aggregation and clusters.} The cluster is the question: all waves, countries, subgroups and option rows of one question in one source. Within a cluster, rows are averaged within each (wave, country) cell with equal weight, and cells are averaged with equal weight. The pooled estimand weights question clusters equally across sources. Because Pew will contribute far more question clusters than Ipsos, the pooled tests H1--H3 are mostly tests on Pew items; we say so in advance and report every quantity per source. Ipsos's month-to-month persistence makes its items hard to win against the last wave. All Pew items in a round share one wave, so a round's Track~A inference covers that round's waves and questions, not future waves. A sensitivity analysis clusters by (question, country).

\paragraph{Effective sample size and noise floor.} Where a margin of error $e$ at 95\% is published, $n^{\mathrm{eff}} = 1.96^2 \cdot 0.25 / e^2$; otherwise the reported $n$ is used and the row is flagged. Ipsos states precision as a Bayesian credibility interval for its online panel ($\pm 3.5$ points at $N = 1{,}000$) \citep{ipsos2026may}, not as a design-based margin of error. We apply the formula to it as the pollster's own statement of precision, which gives $n^{\mathrm{eff}} = 784$ at $N = 1{,}000$; for a non-probability sample this measures neither total error nor a design effect. $n^{\mathrm{eff}}$ enters only the noise floor and the gap variance $v$ below, never a paired $Q$ difference. The noise floor $N_i = \sum_k h_{ik}(1-h_{ik})/(n^{\mathrm{eff}}_i - 1)$ estimates $\mathrm{tr}\,\mathrm{Cov}(h_i)$ under multinomial sampling and is reported beside $Q$. Total poll error is typically larger than reported margins; in US election polls it is about twice as large \citep{shiranimehr2018}. Where two pollsters field the same question in the same country and month, their disagreement is reported as a second, more realistic floor.

\paragraph{Descriptive SimBench score.} For comparability with the literature we report the SimBench score \citep{hu2026simbench} as SimBench defines it: within a dataset $d$ (here a source and wave), row $i$ scores $100\bigl(1 - \TVD(h_i, m_i)/\bar T_d\bigr)$, where $\TVD(p,q) = \tfrac{1}{2}\sum_k |p_k - q_k|$ and $\bar T_d$ is the mean $\TVD(h_j, u_j)$ to the uniform anchor over the dataset's rows, and the score is the mean over rows. Within one dataset this equals the ratio of sums $100\bigl(1 - \sum_i \TVD(h_i,m_i)/\sum_i \TVD(h_i,u_i)\bigr)$. Because $\TVD$ is nonlinear in $h$ and the denominator is itself noisy, the score is biased under sampling noise and is never used for tests. It is computed on our items, not SimBench's, so it is not comparable with SimBench's published figures. On ordinal items we also report the Wasserstein-1 distance on the option index divided by $K_i-1$.

\paragraph{Subgroup fidelity.} Pairs $(g,g')$ are disjoint groups within one demographic dimension (for example two age bands, or men and women) in the same (question, country). On ordinal items $\bar x_g$ is the subgroup mean; on nominal items it is the share choosing the option that is modal in the human country topline; on binary rows it is the share of the first option. Let $d^h$ and $d^m$ be the human and predicted gaps $\bar x_g - \bar x_{g'}$, $\delta = \mathbb{E}\,d^h$, and $v$ the sampling variance of $d^h$. Where the source publishes no subgroup margin, $v$ uses each group's reported $n$ divided by the country's design effect $n/n^{\mathrm{eff}}$, which assumes the country design effect applies to subgroups. Per pair we define
\begin{equation}
G = (d^m - d^h)^2 - v, \qquad D = (d^m)^2 - \bigl((d^h)^2 - v\bigr),
\end{equation}
which are unbiased for $(d^m - \delta)^2$ and $(d^m)^2 - \delta^2$. $G$ is the noise-corrected gap error: it penalizes wrong signs as well as wrong sizes. $D > 0$ on average means the system exaggerates subgroup differences. In a paired comparison of two systems, $v$ cancels from the difference in $G$. Both are averaged over pairs within question clusters and then over clusters, so the bootstrap of Section~\ref{sec:dependence} applies to them directly. As descriptions we also report the gap ratio $\kappa = \sum (d^m)^2 / \sum ((d^h)^2 - v)$, computed only when its denominator exceeds $2\sum v$, with a pairs cluster bootstrap percentile interval that is reported only if fewer than 5\% of replicates have a denominator at or below that bound (the share is reported either way); the attenuation-corrected slope $\sum d^m d^h / \sum ((d^h)^2 - v)$; the share of pairs with the correct gap sign among pairs with $|d^h| > 1.96\sqrt{v}$; and the spread ratio $\mathrm{Var}_m/\mathrm{Var}_h$ on ordinal items.

\subsection{Track B}
For item $i$ with outcome $o_i \in \{0,1\}$ and probability $p_i$, the Brier score is $B(p_i) = (p_i - o_i)^2$ \citep{brier1950verification}, lower is better. The log score is $L(p_i) = o_i \log p_i + (1-o_i)\log(1-p_i)$, higher is better; since every scored probability, imputed and penalty values included, lies in $[0.01, 0.99]$, it is always finite. Both are strictly proper over that range \citep{gneiting2007strictly}. The primary quantity is the cluster-weighted paired difference
\begin{equation}
\bar\Delta = \frac{1}{C}\sum_{c=1}^{C} \frac{1}{n_c}\sum_{i \in c} \bigl(B(p_i) - B(b_i)\bigr)
\end{equation}
against a comparator $b$ over $C$ clusters, negative when the entry beats it. Sampling one contract per event group keeps strike ladders and multi-outcome events from dominating it.

\paragraph{Dependence.}
\label{sec:dependence}
Tests in both tracks use the restricted wild cluster bootstrap-t (null imposed) with Rademacher weights over clusters, as recommended by Cameron, Gelbach and Miller \citep{cameron2008bootstrap}, with 9{,}999 replicates; intervals are obtained by test inversion. Clusters are question clusters in Track~A and rule-assigned event clusters in Track~B. Sensitivity analyses cluster Track~B by resolution week and Track~A by (question, country). The three horizons of one item share its cluster; each confirmatory test uses a single horizon.

\paragraph{Calibration.} CORP reliability diagrams with consistency bands \citep{dimitriadis2021stable} are reported once at least 100 items have resolved for an entry at a horizon.

\subsection{Track C}
The primary score is the ranked probability score \citep{epstein1969scoring} on the three ordered outcome classes; the lift buckets are scored with the same rule over six classes, and secondary-metric directions with the multiclass Brier score. These scores are proper, but they score the partner's reported result, which carries its own sampling error, and they are not corrected for it. As a descriptive companion we report the balanced accuracy (mean per-class recall) of the most probable direction over down, flat and up, which, unlike recall on unfavourable moves alone, cannot be maximized by always predicting deterioration. We also report the share of tests stopped early and the width of partners' intervals.

\section{Leaderboard rules}
\label{sec:rules}

\begin{enumerate}
  \item \textbf{Registration.} An entry registers a hashed system card (Section~\ref{sec:rounds}) before $O_r$. No entry joins mid-round.
  \item \textbf{One ranked system per organization per division per round.} Teams declare their organization on the system card. Ablations, prompt templates and other internal variants, including the organizer's, are unranked; their scores are published beside the ranked entry.
  \item \textbf{Commit before the answer.} All predictions are committed under the slot and salt rules of Track~B before $\tau$.
  \item \label{rule:missing}\textbf{Missing and unrevealed predictions.} An unrevealed or mismatched commitment scores as the worst allowed outcome ($p = 0.99$ on the wrong outcome in Track~B; in Tracks~A and~C, all mass on the option or class farthest from the truth). A missing prediction is imputed as the worse, for the entry, of uniform and the naive last result (Track~A); 0.5 (Track~B); and the worse of uniform and the running prior (Track~C). Abstaining scores exactly what submitting the imputed value would score; under a proper score, submitting one's own belief does at least as well in one's own expectation, so submitting weakly dominates abstaining. An entry is ranked on a track only if it submitted at least 90\% of the items.
  \item \textbf{Divisions and modes are never merged.} Market-blind and market-aware results, and live and retrospective results, appear in separate columns.
  \item \textbf{Minimum count.} An entry is ranked on Track~A after 100 resolved (question, country) cells spanning at least 30 question clusters, on Track~B after 100 resolved items spanning at least 30 clusters, and on Track~C after 100 resolved tests. Track~C will not be ranked in early rounds.
  \item \textbf{Ranks.} At round close, ranks are reported with simultaneous confidence sets for ranks \citep{mogstad2024ranks}, computed from 2{,}000 cluster-bootstrap replicates over the round's common item set, which every entry covers because missing items are imputed. During the round the public display shows point estimates and $n$, labelled interim, with no intervals or ranks.
  \item \textbf{All items, all outcomes.} The full item list, every committed leaf, every reveal, every outcome and every score are published at round close, including items on which the organizer's entries did badly. Track~C outcomes are published under the partner agreement.
  \item \textbf{Organizer entries and recomputation.} Organizer entries are labelled. The scoring code is published, so a third party can recompute every Track~A and~B score from the ledger, given our archive of source files and order-book snapshots; the archived quotes themselves cannot be independently re-fetched. Track~C scores can be recomputed only from partner-reported outcomes.
  \item \textbf{Citing a score.} Any public citation of a benchmark score must name the round, track, division, mode, $n$, number of clusters and the interval.
\end{enumerate}

\section{Pre-registered analysis}
\label{sec:analysis}

\paragraph{Comparison set and intention to treat.} Every comparison uses all non-void items of the round in its stated subset. Missing and unrevealed predictions enter every test, estimate and rank with the values of rule~\ref{rule:missing}, for entries and baselines alike; no item is dropped because a system failed to predict it. All confirmatory tests are one-sided in the stated direction and use the bootstrap of Section~\ref{sec:dependence}.

\paragraph{Error control.} The familywise error rate is 0.05, split in advance between tracks: 0.025 for the five Track~A tests, with Holm's procedure, on the Track~A analysis date, and 0.025 for the two Track~B tests, with Holm's procedure, on the Track~B analysis date. Each track's results can therefore be reported when they exist. A test with fewer than 30 clusters in its comparison set, or below the minimum counts of Section~\ref{sec:rules}, including every test of a system that was not registered, is reported as not tested and enters Holm with $p = 1$, so the family size does not change. Each result is reported with its direction, effect size, a two-sided 95\% interval and the one-sided Holm-adjusted confidence bound, whatever the test outcome.

\paragraph{Confirmatory tests.} Track~A tests use live items and the primary score $Q$, with RW-A denoting the organizer's simulator; Track~B tests use the market-blind division at the 7-day horizon with the Brier score.
\begin{description}
  \item[H1 (A, state and population beyond the same data).] RW-A has lower mean $Q$ than the same-feed LLM.
  \item[H2 (A, strongest cheap rival).] On question clusters with an earlier wave of the same question in the same country whose fieldwork ended within 12 months before $\tau$ (every Pew annual repeat qualifies), RW-A has lower mean $Q$ than the naive last wave.
  \item[H3 (A, persistence).] On the same clusters, the naive last wave has lower mean $Q$ than SP-frontier.
  \item[H4a (A, gap inflation).] SP-frontier exaggerates subgroup gaps: mean $D > 0$.
  \item[H4b (A, gap fidelity).] RW-A has lower mean noise-corrected gap error $G$ than SP-backbone, which shares its model.
  \item[H6 (B, no-market questions).] On calendar items without market coverage, RW-blind beats the template base rate.
  \item[H7 (B, state and population beyond the same feed).] On all items, RW-blind beats the same-feed LLM.
\end{description}
The label H5 is kept for the market bar below, now an estimation target, so that hypothesis numbers match earlier drafts. H4a tests a prediction that prior work makes for individual simulated respondents \citep{chen2026synthetic}; a single prompt that predicts a group distribution is a different construct, so the result is not a replication. H7 tests superiority only. The equivalence test for the opposite direction belongs to the matched comparison in \emph{Don't Ask a Model, Ask a World}, which holds the question set and models fixed in a way an open benchmark cannot.

\paragraph{What each test decides.} If H1 or H7 fails, the simulator's state and population have not been shown to add anything beyond giving the same model the same data and the same number of calls. If H2 fails, the simulator has not been shown to beat assuming nothing changed, the cheapest rival there is; H3 checks that this rival is as strong here as it is reported to be elsewhere. H4a and H4b say whether the gap exaggeration found for prompted models appears in group-level prediction, and whether the simulator avoids it. H6 says whether the simulator adds value where no market exists, which is where a simulator would be most useful.

\paragraph{Estimation targets, no test.} Reported with 95\% intervals: H5, the market against SP-frontier on market-covered items, $\bar\Delta(\text{SP-frontier} - \text{market})$, which is expected to favour the market and is estimated rather than tested to save error rate for the simulator hypotheses; RW-blind against the market at 7 days; RW-aware against the market in the market-aware division (the confirmatory test of the organizer's market-blind record against the market is H6 of \emph{Receipts}, on its own question set); and all Track~C quantities. A null result is not evidence of equality; we report the interval and the minimum detectable effect.

\paragraph{Exploratory.} Other horizons and pooled horizons, log scores, the naive status-quo baseline for calendar items, RW-A against SP-backbone, the non-Western country subset of H1 \citep{durmus2023global}, SP-frontier against SP-backbone, templates T2 and T3, the drift baseline, the demographic lookup, the spread at a 0.10 coverage threshold, the last-15-minute market reference, $\kappa$ and the gap slope, and the ablations below. Exploratory results are labelled as such and carry no claims.

\paragraph{Ablations (exploratory, unranked).} \emph{No assimilation}: the segment positions are initialized from source data at $O_r$ and never corrected against polls. \emph{No within-segment spread}: every agent offset is set to zero, so predictions use segment positions alone and the agent records contribute nothing. H7 and these ablations are the benchmark counterparts of the matched comparison in \emph{Don't Ask a Model, Ask a World}.

\paragraph{Hypothesis blindness.} As of this draft a first version of the organizer's simulator exists in code but has forecast no live question, and no predictor, baseline or otherwise, has been scored by us on any Track~A or~B item, including the Pew spring 2026 development set. Its only results are two synthetic checks of its assimilation operator, which use no poll, survey or market data. While writing the item and void rules we read published Pew spring 2026 reports and the Ipsos May 2026 report; they shaped those rules, not the direction of any hypothesis. H3 and H4a take their direction from published prior work. The organizer plans to start a hashed daily forecast record in mid-October 2026; its results will exist before round~1 opens and may inform the system's design, which is why this text is to be frozen before that record's first resolved question (Section~\ref{sec:rounds}).

\paragraph{Track C disclosure.} An earlier draft contained a confirmatory Track~C hypothesis. It was written after an internal prototype had been scored on the retrospective pilot set, and its form followed that internal result, so it was not formed blind. It has been removed; round~1 reports Track~C only as estimates.

\paragraph{Minimum detectable effects.} With power 0.8 at Holm's strictest level in each track, one-sided $0.025/5$ in Track~A and $0.025/2$ in Track~B, the minimum detectable mean difference is $\mathrm{MDE} = c\,\sigma\sqrt{D/n}$ with $c = 3.42$ (Track~A) or $3.08$ (Track~B), where $\sigma$ is the standard deviation of the per-cluster difference, $n$ the number of clusters and $D$ the remaining design effect. With $D=1$, the Track~A MDE is $0.62\sigma$, $0.34\sigma$, $0.20\sigma$ and $0.11\sigma$ at 30, 100, 300 and 1{,}000 question clusters; Track~A will have tens of question clusters per round, not hundreds, so only large effects are detectable. For Track~B, if, purely as a planning assumption, $\sigma = 0.1$ in Brier units, the MDE is 0.022, 0.014 and 0.010 at 200, 500 and 1{,}000 independent clusters. Detecting a gap of the size Prophet Arena reports between its best model and the market, 0.003 \citep{yang2025prophet}, would need about 10{,}600 independent clusters at this level, and about 8{,}700 at an unadjusted two-sided 0.05, far beyond one round. These figures use a normal approximation; with 30 to 100 clusters the bootstrap-t is somewhat more conservative. A round-1 null on H6 or H7 therefore means only that no difference larger than the MDE was found. Planning values for $\sigma$ are not results. Before submission we will estimate the per-cluster standard deviation of $Q$ differences between existing baselines (last wave, SP-frontier, uniform) on the Pew 2026 development set, and count how many calendar items over January--June 2026 lacked market coverage, since H6 may never reach 30 clusters; both are pending artifacts in Table~\ref{tab:pending}. Any change to planned counts goes into the deviations log.

\begin{table}[t]
\centering
\small
\begin{tabular}{@{}L{5.6cm}L{8.6cm}@{}}
\toprule
\textbf{Result} & \textbf{Status and producing artifact} \\
\midrule
Feasibility: Pew 2025--2026 repeated questions, country turnover, $\sigma$ of baseline $Q$ differences & Pending (before submission): \nolinkurl{bench/feasibility/track-a-pew-2025-2026.json} \\
Feasibility: calendar items without market coverage, Jan--Jun 2026 & Pending (before submission): \nolinkurl{bench/feasibility/track-b-calendar-coverage-2026h1.json} \\
Round-1 manifest and system cards & Pending: \nolinkurl{bench/rounds/round-01/manifest.json} \\
Item census per track (items, clusters, subgroup clusters, void counts by rule and topic) & Pending: \nolinkurl{bench/track-a/census/round-01.json}, \nolinkurl{bench/track-b/census/round-01.json} \\
Track A $Q$, noise floors, two-pollster floor, SimBench score (live and retro), per source and region & Pending: \nolinkurl{bench/track-a/results/round-01/scores.json} \\
Track A subgroup fidelity ($G$, $D$, $\kappa$, slope, sign agreement, spread) & Pending: \nolinkurl{bench/track-a/results/round-01/subgroups.json} \\
Track B commit log, batch roots, time-stamp tokens, market snapshots & Pending: \nolinkurl{bench/track-b/ledger/} \\
Track B $\bar\Delta$, Brier, log score per horizon, division and region & Pending: \nolinkurl{bench/track-b/results/round-01/scores.json} \\
Track B calibration (after 100 resolved) & Pending: \nolinkurl{bench/track-b/results/round-01/calibration.json} \\
Confirmatory tests (seven, Holm within track) & Pending: \nolinkurl{bench/analysis/round-01/confirmatory.json} \\
Track C retrospective scores (pre- and post-cutoff) & Pending: \nolinkurl{bench/track-c/results/retro-v1/scores.json} \\
Track C live ledger and estimates & Pending: \nolinkurl{bench/track-c/ledger/}, \nolinkurl{bench/track-c/results/live-01/scores.json} \\
Single-prompt template spread & Pending: \nolinkurl{bench/baselines/prompt-spread/round-01.json} \\
Simulator ablations & Pending: \nolinkurl{bench/ablations/round-01/scores.json} \\
Deviations log & Pending: \nolinkurl{bench/AMENDMENTS.md} \\
\bottomrule
\end{tabular}
\caption{Every result this paper will report, and the committed artifact that will produce it. No number in this table exists yet. Every interval in these artifacts is reported with its number of entries, items and clusters.}
\label{tab:pending}
\end{table}

\subsection{Results}
\pending{the artifacts listed in Table~\ref{tab:pending}}
No score, ranking or comparison is reported in this draft. When a track reaches its analysis date, this section will report every pre-registered test, hit or miss, with the commit hash of the artifact that produced it.

\section{Dataset release}

Each track is released with a datasheet \citep{gebru2021datasheets} and Croissant metadata \citep{akhtar2024croissant}. Releases are versioned by round; a published item is never edited, and corrections are new rows that reference the old one. The Track~A scorer will be packaged as an Inspect task \citep{inspect2026}. The scoring code is planned to depend only on the Python standard library, so the ledger can be verified without our infrastructure; it has not been written yet.

\section{Limitations and failure modes}
\label{sec:limits}

\paragraph{Slow accumulation and low power.} Forward-dated items resolve at the speed of the world. One round will hold hundreds, not thousands, of resolved Track~B clusters, tens of Track~A question clusters, and few live Track~C tests. The minimum detectable effects above are large relative to the gaps reported between strong models and markets.

\paragraph{Live Track A may be thin and Pew-dominated.} Pollsters change questions between waves, Pew's country set changes each year, and subgroup tables are published for one round-1 source only. The census may leave too few question clusters to test H1--H4b, and the tests that do run are mostly about Pew items in one wave.

\paragraph{The benchmark does not validate forked scenarios.} The product's central feature, forking today's world with a hypothetical event and running an ensemble over seven simulated days, produces predictions about counterfactuals that never occur. They have no ground truth, and nothing here validates them. Benchmark scores on factual forecasts do not transfer to counterfactual claims.

\paragraph{Self-preferencing and Goodhart effects.} The organizer chose the tracks, metrics (including the subgroup gap scores), divisions and item rules, and these may suit its own design. Rule-based item generation, an external beacon and a hashed manifest limit choice within a round, not the choice of the rules. The operator and the simulator team belong to one company, so the information barrier rests on access controls we run ourselves. An independent host is a goal, not yet a fact.

\paragraph{Market reference.} The market midpoint is not a probability: prices embed fees and need not equal traders' mean belief \citep{manski2006interpreting}, and prices near 0 or 1 show a favourite-longshot pattern \citep{snowberg2010,page2013calibrated}. Thin markets can be moved near $\tau$; the time-weighted window and volume floor reduce but do not remove this. Archived snapshots are hashed, but third parties cannot re-fetch historical order books, so they must trust our archive for the quotes themselves. Market coverage is narrow and mostly US-centred.

\paragraph{Leakage to market-blind entrants.} RW-blind drops market feeds and news mentioning markets or odds, but news can reflect market moves without naming them. Item generation also leaks: a market item exists because a venue listed it, and a strike ladder's listed strikes cluster around the market's expected value, so a sampled strike tells a market-blind entrant roughly where the market's centre lies. Market-blind means blind by input, not by information.

\paragraph{Truth is noisy.} Polls carry sampling error, design effects and house effects, and total error exceeds stated margins \citep{shiranimehr2018}. $Q$ differences are unbiased for differences against the poll's expectation under sampling noise only; house effects and wording changes are not corrected, and renormalization and rounding add small bounded errors.

\paragraph{Timestamps rest on authorities.} Minute-level commit times rest on at least one of two RFC~3161 authorities being honest. The OpenTimestamps attestation needs no trust but bounds the time only to hours.

\paragraph{Contamination is reduced, not eliminated.} Declared cutoffs can differ from effective ones \citep{cheng2024dated}. Retro-mode freezes cannot be fully audited, and the retrospective Track~C set is almost certainly in some training data. Only live items enter tests and ranks.

\paragraph{Hashes prove integrity, not foresight.} A timestamped hash proves that one prediction per slot existed at a time. It does not prove that the prediction used no improper information, and it does not stop an entrant from developing many private systems between rounds and registering the best, or one organization from registering several nominal teams; the one-per-organization rule relies on honest declaration.

\paragraph{Track C depends on partners.} Partners, and the tests they share, are not a random sample. Outcomes are self-reported and unverifiable. A partner could stop a test early for reasons correlated with its result; registered end dates make that visible, not impossible.

\paragraph{Licensing.} Redistribution terms for market data, Pew data, Ipsos data and other sources are not settled. Where they are unclear we release links and hashes, which makes independent replication harder.

\paragraph{Not yet a dataset.} This draft releases no items, no baseline scores and no benchmark code; the organizer's engine repository has a same-feed LLM baseline for its own daily record, but not the benchmark's baselines, item generator or scorer. We will submit to a datasets and benchmarks track only after releasing, as committed artifacts, the retrospective Track~A development set (Pew spring 2026) with baseline scores and the two feasibility artifacts of Table~\ref{tab:pending}. If they are not ready, we will submit this protocol as a Stage-1 registered report instead.

\section{Ethics statement}

\paragraph{Personal data.} The benchmark releases only published aggregate estimates and blinded experiment summaries. The demographic-lookup baseline is fit privately on licensed respondent-level files where the license allows; those files are never redistributed, and released cell tables have at least 50 respondents per cell. No user-level product data is stored.

\paragraph{Influence on elections and markets.} Track~B includes elections. Forecasts inside the benchmark are published only as hashes until resolution. This does not constrain the organizer's public product, which may publish election forecasts outside the benchmark; those forecasts are governed by the product's own policy, discussed in the \emph{Governing Public World Simulators} companion. The benchmark does not require or reward trading, and the organizer's entries do not trade. It cannot prevent external entrants from trading, so each entry must disclose on its system card any positions its organization holds in markets on benchmark items.

\paragraph{Subgroup prediction is dual-use.} A system that predicts the opinions of narrow subgroups well could be used to target persuasion. Track~A scores only broad published subgroups with at least 100 respondents, never individuals. Track~C scores experiments at the level of an arm.

\paragraph{Representation.} A pooled leaderboard can reward systems that are accurate for well-polled populations and poor elsewhere. Region-level scores are reported beside every pooled score (Table~\ref{tab:pending}).

\paragraph{Overclaiming.} A benchmark gives systems a number to advertise. Any public citation of a score must carry the details in rule~10 of Section~\ref{sec:rules}. A good score on factual forecasts must not be presented as validation of counterfactual ``what-if'' simulations, which this benchmark does not test.

\section{Competing interests and funding}

The benchmark is designed, funded and operated by Real World, the company that builds the simulator it evaluates, and the authors work for that company. The organizer chose the tracks, metrics, item rules and analysis plan, runs the operator, and enters its own system; Track~C partners may be or become its customers. No external funding supported this work. The safeguards in Sections~\ref{sec:rounds}--\ref{sec:analysis} limit what the organizer can choose after the rules are frozen; they do not remove the conflict, which readers should weigh.

\section{Conclusion}

Published evidence for population simulators comes mostly from questions whose answers were public before the systems were scored. The Real World Benchmark asks them to predict polls whose fieldwork has not begun, events that have not happened and experiments that have not started, and to do better than cheap rivals: nothing changes, one prompt to the same model, the same model given the same data, and the market. This text fixes the tracks, item rules, scores, tests and analysis dates before any system has been scored, and lists every parameter still open with the point at which it is frozen. If the simulators cannot beat these baselines, that is a result the field needs to know.

\bibliographystyle{plainnat}
\bibliography{refs}

\end{document}
